% Short running form; the full title is too wide for the head's right box.
\runninghead{The Feast That Now Is}
\properlane{feast-that-now-is}{The Feast That Now Is: the Church at the Lord's Table and Its Evening Sacrifice}
\label{sec:feast-that-now-is}

Heard from its middle, the Gospel is about a hall that is full. The servants
have brought in everyone they found, bad and good, the guests are at table, and
the king comes in to see them. St Augustine and St Gregory the Great both say
that this wedding is the Church now, with good and bad at one table, and both
say so to congregations gathered in church: Augustine calls the feast the
Lord's Table and the Feast of the Holy Scriptures, Gregory the banquet of the
Scriptures and of the divine Word, and St John Chrysostom addresses its guests
as those who have partaken of the mysteries. Read from there, the chants and
prayers of this Mass are the voice of that assembly at that table. It cries in
its tribulation and is heard; its prayer rises like incense and the lifting up
of its hands joins the evening sacrifice, which St Augustine reads as the
Lord's Passion; it offers its gifts before the eyes of God's majesty and asks
that they be saving; and at the Communion it asks to keep what it is
commanded, and to be freed by the sacrament's healing work. The garment is
the question put to every guest before he approaches. St Augustine, St Gregory
and St John Chrysostom carry this reading from their preaching on the Gospel
and the psalms. Bl.\ Ildefonso Schuster, commenting on this Mass as a Mass,
reads its Gradual, Secret and Postcommunion in the same direction, though the
banquet of its Gospel is for him the heavenly one.

\subsection{Two feasts: St Augustine and St Gregory on the present Church}

St Augustine begins his sermon on the parable with the table. ``All the
faithful know the marriage of the king's son, and his feast, and the spreading
of the Lord's Table is open to them all who will. But it is of importance to
each one to see how he approaches, even when he is not forbidden to approach
It. For the Holy Scriptures teach us that there are two feasts of the Lord; one
to which the good and evil come, the other to which the evil come not''
(\work{Sermo} 90, 1). The feast of the Gospel just read is the first of these,
and Augustine addresses those at it in two directions at once: ``that ye look
not for the good without, that ye bear with the evil within''. Its guests are
already at the Lord's Table ``which is in this life'' (90, 5), and the whole
purpose of the sermon is gathered into one sentence: ``Be that feast which now
is, received worthily, that we may attain to the other'' (90, 4).

St Gregory says the same of the time of the wedding. Luke's supper, he says, is
the eternal and last banquet; Matthew's wedding is the present Church, because
from it some who have come in go out again (\work{Homiliae in Evangelia}
38, 1). The hall full of bad and good shows it: \latin{per has regis nuptias
praesens Ecclesia designatur, in qua cum bonis et mali conveniunt}, by this
king's wedding the present Church is signified, in which the bad come together
with the good. The good are not to leave the bad behind while the feast lasts.
\latin{Si ergo boni estis, quandiu in hac vita subsistitis, aequanimiter
tolerate malos}: if you are good, bear the bad with an even mind for as long
as you remain in this life, for the grain is threshed under the chaff and the
rose grows with the thorn (38, 7). Then Gregory turns to the people in front
of him. \latin{Sed quia jam, largiente Domino, nuptiarum domum, id est sanctam
Ecclesiam intrastis, solerter, fratres, aspicite, ne aliquid de mentis vestrae
habitu rex ingrediens reprehendat}: since by the Lord's gift you have already
entered the house of the wedding, that is, holy Church, look carefully,
brethren, lest the king coming in find fault with anything in the dress of
your mind (38, 9). And he names the feast at which they sit. \latin{Nos sumus,
fratres charissimi, qui in nuptiis Verbi discumbimus, qui jam fidem in
Ecclesia habemus, qui Scripturae sacrae epulis pascimur}: we are those who
recline at the wedding of the Word, who already have faith in the Church, who
are fed with the banquet of Holy Scripture; let each weigh his heart, whether
he holds hatred against anyone, whether envy of another's happiness burns in
him (38, 11). In a later sentence he says it in four words, \latin{divini Verbi
epulas sumimus}, we take the banquet of the divine Word (38, 14).

St John Chrysostom, closing his homily on the parable, speaks to his hearers as
guests at the same wedding, and he names what they have received: ``Hear ye,
as many as having partaken of the mysteries, and having been present at the
marriage, clothe your souls with filthy deeds. Hear whence ye were called. From
the highway. Being what? Lame and halt in soul \dots\ Reverence the love of
Him, who called you, and let no one continue to have filthy garments''
(\work{Homilies on Matthew} 69). Augustine, in his other sermon on the parable,
speaks of the call in the same present tense, and puts himself in the
servants' place: ``seeing ye have been invited through me; for though He
invited you, He invited you by my ministry. Ye are all at the feast, have the
wedding garment'' (\work{Sermo} 95, 5). The Bridegroom, he adds in the same
sermon, is the one ``who calleth together to the feast, and quickeneth whom He
calleth'' (95, 6).

The three preachers do not name the feast alike. Augustine's is the Lord's
Table and the feast of the Scriptures, Gregory's the banquet of the Word and of
Holy Scripture, Chrysostom's the marriage at which his hearers have received
the mysteries. What they share is the time: the feast is now, its guests are
in the Church, and the king has not yet come in.

Bl.\ Ildefonso Schuster, who hears in this Mass's Gradual the Christian's
evening oblation and in its Secret and Postcommunion the Eucharist, places the
Gospel's feast elsewhere. ``The primary aim
of the predestination of souls to the heavenly banquet'', he writes of the
parable, ``is the supreme glorification of Christ as first-born of the human
family and head of the Church''; and of the guest without the garment, ``God
calls our souls, but they must be in accord with their calling, so that the
grace of eternal happiness may also be their due reward'' (p.~173). His banquet
is the one to come, the feast that St Augustine sets beside the present one as
the other, to which the evil do not come. It is Augustine and Gregory, not
Schuster, who find the Church now in this Gospel's wedding.

\subsection{The evening sacrifice: St Augustine on the Gradual}

The Gradual is the assembly's evening prayer: ``Let my prayer be directed as
incense in thy sight; the lifting up of my hands, as evening sacrifice.'' St
Augustine hears in the psalm the voice of the whole Christ. ``This not I alone
say: whole Christ saith it. But it is said rather in the name of the Body'',
and the Body must go on praying as long as its trouble lasts: ``if tribulation
remain for the Church, for the Body of Christ, even to the end of the world,
let it not only say, `I have cried unto Thee,' but also, `Listen unto the voice
of my prayer'\,'' (\work{Enarrationes} 140, 2).

On the verse the Gradual sings, Augustine turns from the Body to the Head. ``That
this is wont to be understood of the Head Himself, every Christian
acknowledgeth. For when the day was now sinking towards evening, the Lord upon
the Cross `laid down His life to take it again'\,''. The members are not left
out: ``Still we too are figured there. For what of Him hung upon the tree, save
what He took of us?'', and on the Cross ``our old man is crucified with Him''.
Then he names the sacrifice: ``That then is the `evening sacrifice,' the
Passion of the Lord, the Cross of the Lord, the offering of a salutary Victim,
the whole burnt offering acceptable to God. That `evening sacrifice' produced,
in His Resurrection, a morning offering. Prayer then, purely directed from a
faithful heart, riseth like incense from a hallowed altar'' (140, 3). The
evening sacrifice is Christ's; the prayer that rises with it is the Church's;
and the altar from which it rises is the hallowed altar of a faithful heart.

St Robert Bellarmine, expounding the same verse, finds the high priest of that
prayer: ``Christ is the high priest, for it is through him, as being our
advocate, that we must always pray; and it is for that reason that we conclude
every prayer with `through our Lord Jesus Christ'\,''. The prayer rises
straight only when the intention is pure and the attention constant, since
distractions ``blow away, and dissipate the incense of their prayer''. And he
offers, as one possibility among others, the figure Augustine states outright:
the psalmist asked that his prayer be like the evening sacrifice ``perhaps,
because the evening sacrifice was of more value as being a figure of the
sacrifice on the cross, that occurred in the evening'' (on Ps 140:2). St John
Chrysostom, expounding the psalm in Greek, adds why the Church sings it: the
Fathers appointed it to be said every evening not for its one phrase about the
evening sacrifice, but as a saving medicine and a cleansing of the sins of the
day (\work{Expositio in Psalmum} 140). Schuster, reading this Mass's Gradual,
calls every Christian offering evening: ``In this present life our oblation to
God is always an evening sacrifice, for it is enveloped in the twilight of
faith'', and ``before attaining to the morning joy of the beatific vision, it
is fitting that we should endure with Christ the labour and the sorrow of the
evening hours of the preceding Parasceve'' (p.~172).

In a Mass celebrated with incense the Gradual's verse is heard again: at the
Offertory the priest says it, with the two verses that follow, as he incenses
the altar (\work{Ordo Missae}, no.~1037). At such a Mass the evening prayer
that the assembly sang between the lessons is said again by the priest over the
altar of the sacrifice, and the incense of the psalm rises in fact.

St Hilary of Poitiers answers the question of the evening sacrifice otherwise. For him
the evening is the last age of the world, and the sacrifice that pleases God is
the works of mercy of the last judgement, done with lifted hands. The two
answers are not opposed in doctrine: Augustine reads the verse first of the
Head, and Hilary of the hands of the members. But they are two answers to what
the verse's sacrifice is, and the feast that now is takes its evening
sacrifice from Augustine's.

\subsection{Tribulation and gifts: the Offertory and the Secret}

The Offertory opens the offering of the gifts, and it takes up the Introit's
word: \latin{Si ambulavero in medio tribulationis, vivificabis me}. St
Augustine's exposition of the verse reads like a gloss on the Introit's
promise: ``whatsoever tribulation thou art in, confess, call on Him; He freeth
thee, He reviveth thee'', and he widens the tribulation
to the whole of this life: ``Love the other life, and thou shalt see that this
life is tribulation, whatever prosperity it shine with, whatever delights it
abound and overflow with'' (\work{Enarrationes} 137, 10). The antiphon at the
Mass's entrance promised that the Lord would hear \latin{de quacumque
tribulatione}; the chant at the offering is sung by those who are still
walking in it. St John Chrysostom notices that the psalmist does not say that
God will drive the tribulation away, but that he will keep him alive while he
remains in the midst of its dangers, and that this is the wonder (\work{Expositio
in Psalmum} 137). St Robert Bellarmine names the right hand that saves:
``your strength and power, Christ, hath saved me'' (on Ps 137:7).

The Secret then offers the gifts \latin{oculis tuae maiestatis}, to the eyes of
God's majesty, the same eyes before which the Gradual asked the prayer to rise
\latin{in conspectu tuo}. It asks one thing, that they be \latin{salutaria
nobis}. Schuster reads the prayer of the Eucharist and of those who share it:
``we pray that the holy oblation may not only bring glory to God, but may also
be a pledge of eternal salvation to all those that share therein'', and since
the sacrament's effect in each person answers to his dispositions, ``we here
ask for such dispositions as may enable the mystical sacramental oblation to
exert without hindrance its full power in our souls'' (p.~173). The
\latin{Placeat} at the end of every Mass asks in the same words that the
sacrifice offered \latin{oculis tuae maiestatis} be propitiable for all for
whom it was offered (\work{Ordo Missae}, no.~1137). The Secret's
\latin{salutaria nobis} is Augustine's ``received worthily'' turned into
petition: the gifts are saving to those who come to the table in the garment.

\subsection{Commanded and asked, healed and held: the Communion and the Postcommunion}

The Communion is sung after the Communion itself, and its two verses are a
command and a wish: ``Thou hast commanded thy commandments to be kept most
diligently. O! that my ways may be directed to keep thy justifications.'' Its
\latin{dirigantur} answers the Gradual's \latin{dirigatur}. The prayer that was
to rise straight as incense before the sacrifice becomes, after the Communion,
the prayer that the ways of those who have eaten be set straight. St Augustine
reads the wish as the words of one who knows that he cannot do what he is
commanded without help, and must pray God to grant what he enjoins
(\work{Enarrationes} 118, Aleph 5), and St Hilary says on the same verse that
we must be helped and directed by God's grace. Schuster hears joy in the wish.
Those who have been at the table do not take the direction of their ways for
granted; they ask for it.

The Postcommunion names the sacrament's work in the language of medicine,
\latin{medicinalis operatio}, and asks that it free us \latin{a nostris
perversitatibus} and make us cleave to God's commandments. Schuster reads it of
the Eucharist: ``we make our petition that the healing grace of the Eucharist
may mercifully deliver us from our perversity and make us ever to be devoted
to the divine will'' (p.~174). The Mass that began in tribulation ends with a
prayer for healing, and the healing asked is the power to keep what is
commanded.

\subsection{The other elements at the table}

The Introit's promise, \latin{de quacumque tribulatione clamaverint ad me,
exaudiam eos}, is the promise on which the assembly's cry rests, and Schuster
reads it as the Lord's assurance that no one who calls on him is disappointed
(p.~171). The Collect asks that body and soul be set free together,
\latin{mente et corpore pariter}, for the service of God; Schuster reads the
prayer as a petition for the order in which the whole man is directed to God
as his final end without hindrance (pp.~171--172). At the table that freedom
is the freedom to approach.

The Epistle's commands are the manners of the table. ``For we are members one
of another'': Chrysostom's image of the body whose eye will not lie to its foot,
and Aquinas's \latin{Membra enim se invicem diligunt}, describe guests who
belong to one body. ``Let not the sun go down upon your anger'': St Augustine
tells his hearers what not to say as they come to the table, ``as thou art
about to approach the Table of the Lord, the Feast of the Holy Scriptures, do
not say, `O that mine enemy might die! Lord, if I have deserved ought of Thee,
slay mine enemy'\,'', and he sets before them the Bridegroom ``hanging upon the
Cross for thee, and praying to His Father for His enemies'' (\work{Sermo} 90,
9). Chrysostom, expounding the Gradual's psalm, will not let hands be lifted in
prayer with anger in them. And ``Give not place to the
devil'': St Jerome and St Thomas Aquinas, each commenting on that verse,
recall the one guest at the Lord's own supper to whom the devil was given a
place. Jerome places it \latin{in convivio Salvatoris}, at the Saviour's
banquet (on Eph 4:27, PL 26, col.~512), and Aquinas cites the Gospel's
\latin{post buccellam introivit in eum satanas}, after the morsel Satan entered
into him (\work{In Epistolam ad Ephesios} c.~4, lect.~8; Jn 13:27).

The Alleluia gives the order of the assembly's prayer. St Augustine observes on
its verse that the psalm begins with praise and then calls upon God: ``For
after commencing with praise, calling upon God is wont to follow, whereunto he
that prayeth doth next add his longings: whence the Lord's Prayer itself hath
at the commencement a very brief praise, in these words, `Our Father which art
in Heaven'\,'' (\work{Enarrationes} 104, 1). \latin{Confitemini \dots\ et
invocate}: the Alleluia, sung before the Gospel, gives the order that the
prayers of the Mass keep.

\subsection{The table is not the garment}

St Augustine, who calls the present feast the Lord's Table, denies in the same
sermon that the garment is the Altar or what is received at it, since many eat
and drink judgement to themselves (\work{Sermo} 90, 5). In another sermon he
says it again: ``The Sacraments of the Altar the good and bad receive
together'' (\work{Sermo} 95, 7). Receiving at the table does not clothe the
guest. The feast now includes the table, and the garment is what the guest
brings to it; that is why St Gregory warns those who are already inside to
look carefully lest the king, coming in, find fault with the dress of their
minds (38, 9).

The witnesses are also exact about what they call the feast. Augustine names
the Lord's Table and the Scriptures, Gregory the Word and the Scriptures, and
Chrysostom the mysteries of which his hearers have partaken. None of them says
in so many words that the wedding feast of this Gospel is the Eucharist, and
Gregory, whose feast is the banquet of the Word and of Scripture, does not
name the altar at all. The table at which the assembly sits is the Church's
whole feast of Word and sacrament, of which those three preachers each name a
part.

Finally, the incense. The Gradual's verse returns in the Order of Mass only
where the Mass is celebrated with incense. At a low Mass the psalm's incense is
the prayer itself, as Augustine reads it, rising ``from a hallowed altar'' of
a faithful heart. St Hilary, for his part, places the wedding in the
resurrection and the king's entry in the assembly of the risen, and the present
feast looks toward that one.

\subsection{The four senses of this reading}

\begin{fourSenses}
\item[Literal.] The king's hall is filled with guests good and bad, and the
king comes in to see them. The psalmist asks that his evening prayer rise
before God like incense and that the lifting up of his hands be accepted as the
evening sacrifice; walking in the midst of tribulation, he trusts to be kept
alive by God's hand and saved by his right hand.

\item[Allegorical.] The feast is the Church now, at the Lord's Table and the
banquet of the Scriptures and of the Word (Augustine; Gregory), whose guests
have partaken of the mysteries (Chrysostom). The evening sacrifice is the
Lord's Passion, the offering of a salutary Victim, in which the members too
are figured (Augustine), perhaps prefigured in the evening offering of the
temple (Bellarmine); the prayer of the whole Christ rises with it, offered
through Christ the high priest (Bellarmine).

\item[Moral.] To examine the dress of the mind before the king comes in
(Gregory), and to come to the table without hatred or envy (Gregory) and
without praying for an enemy's death (Augustine); to bear the bad within the
one Church with an even mind (Gregory; Augustine); to lift hands in prayer
without anger (Chrysostom on Ps~140) and to pray with a pure intention that
lets the incense rise straight (Bellarmine); and in every tribulation to
confess and call upon God (Augustine on Ps~137).

\item[Anagogical.] The other feast, to which the evil do not come, is the one
the feast that now is prepares, received worthily so that the guests may
attain to it (Augustine, \work{Sermo} 90). The evening sacrifice of the Passion
produced in the Resurrection a morning offering (Augustine on Ps~140), and the
evening in which every present oblation is made gives way to the morning of
the beatific vision (Schuster).
\end{fourSenses}
